\section{Introduction}
\label{sec:intro}

The capacity of modern agents depends not only on the underlying model~\cite{deepseekai2026deepseekv4,glm5team2026glm5vibecodingagentic, yang2025qwen3,team2023gemini}, but on the mediation imposed by the surrounding \textit{harness}~\cite{lu2026openclaw,liagent,claudecode}. This harness converts raw model outputs into structured agent behaviors by determining how tasks are represented, how external services are accessed, and how intermediate decisions are communicated during execution. As agents tackle longer-horizon tasks in richer environments, harness design becomes integral to agent development.

Despite this importance, harness development remains far from a mature engineering discipline. \textit{First}, harnesses are hand-engineered and static: any change in model version, tooling, or problem domain requires bespoke modification, with no mechanism for experience-driven improvement. \textit{Second}, harnesses are architecturally entangled: they typically combine prompt templates, tool wrappers, retry policies, and memory in the same codepaths, so changes to one component silently break others, and reuse across domains reduces to copying rather than composition. \textit{Third}, harness engineering and model training operate independently: trajectory data collected while improving the harness is discarded rather than incorporated into model training, and model improvements do not translate into harness improvements.

We address these gaps by treating the harness as a \textit{first-class object} that can be composed, adapted, and evolved alongside the model. \ProjectName embodies this principle as a unified harness foundry. It begins with a modular foundation: harness primitives spanning context, tools~\cite{feng2025retool}, skills, control, and memory are described via typed interfaces and composed via a substitution algebra. This separates concerns that existing systems typically conflate. On top of this substrate, we introduce AEGIS, an observability-driven and auditable harness adaptation engine. Framing harness adaptation not as ad-hoc editing but as a learning problem over symbolic artifacts (prompts~\cite{zhou2025proposer}, tools, memory, and control policies) reveals that standard RL pathologies (reward hacking, catastrophic forgetting~\cite{kirkpatrick2017overcoming}, under-exploration~\cite{ladosz2022exploration}) become concrete design risks. To address these risks, AEGIS combines full trace observability with a four-stage pipeline (Digester, Planner, Evolver, and Critic) that compresses traces, plans adaptations, generates candidates, and assesses changes. Finally, we close the loop between harness adaptation and model training via \textbf{harness-model co-evolution}. Traces produced during harness adaptation serve as reinforcement-learning signal for model training, so that model improvements feed back into subsequent harness evolution.

We empirically validate \textbf{\ProjectName} across five benchmarks (GAIA, ALFWorld, WebShop, $\tau^3$-Bench, SWE-bench Verified), three task-agent families (Claude Sonnet 4.6, GPT-5.4, Qwen3.5-9B), and up to 15 evolution rounds. Harness evolution yields an average absolute gain of +14.5\% across 15 model--benchmark configurations, with individual gains ranging from 0.0\% to +44.0\% among improving configurations (14 of 15), from +1.1\% ($\tau^3$-Bench, near-ceiling baseline) to +44.0\% (ALFWorld, weakest agent). Gains exhibit an inverse-scaling pattern: on ALFWorld and GAIA, the weakest task agent benefits most (+44.0\% for Qwen3.5-9B vs.\ +11.2\% for Sonnet 4.6 on ALFWorld), suggesting that evolved harnesses address behavioral gaps that weaker models cannot self-correct. On heterogeneous task sets (GAIA), single-harness evolution stagnates; a variant-isolation ablation restores stable improvement (+13.6\%, non-degrading over 15 rounds).


In summary, our contributions are four-fold:
\begin{itemize}
  \item \textbf{Harness Composition} (Section~\ref{sec:method-foundry}). We formalize the harness as a first-class, typed object composed of processors attached to lifecycle hooks. A nine-dimensional taxonomy spans the full behavioral space, and a substitution algebra enables per-task configuration with type-safe insertion and removal. This compositional structure makes the intended scope of each behavioral change explicit---a precondition for the variant isolation that stabilizes evolution.
  \item \textbf{Harness Adaptation} (Section~\ref{sec:method-aegis}). We introduce AEGIS, a trace-driven, multi-agent harness evolution engine. An operational mirror maps harness adaptation onto standard RL constructs, converting familiar RL pathologies (reward hacking, catastrophic forgetting, under-exploration) into concrete design risks addressed by a four-stage pipeline (Digester, Planner, Evolver, Critic) with deterministic gating. An optional variant-isolation strategy prevents cross-task interference on heterogeneous benchmarks.
  \item \textbf{Harness-Model Co-Evolution} (Section~\ref{sec:method-coevolution}). We close the optimization loop by interleaving harness evolution with model reinforcement learning over a shared replay buffer. Cross-harness GRPO enables the model to internalize strategies from successive harness versions, breaking the scaffolding ceiling that limits harness-only adaptation and the training-signal ceiling that limits model-only RL.
  \item \textbf{Empirical Validation} (Section~\ref{sec:experiments}). Across five benchmarks, three task-agent families, and up to 15 evolution rounds, \ProjectName yields an average gain of +14.5\% (up to +44.0\%), with gains largest where baselines are lowest. A variant-isolation ablation resolves stagnation on heterogeneous task sets, and co-evolution yields an additional +4.7\% over harness-only evolution (Section~\ref{sec:exp-coevolution}).
\end{itemize}